\section{Finite complex waves and robust nodal-area witnesses}
\label{sec:waves}

This section converts a fixed strict profile into a real approximate
eigenfunction.  Its weighted residual can be made smaller than any prescribed inverse
power of the frequency, while its value and gradient cannot vanish together.
The same function has a nodal-area witness whose constant is independent
of the size and derivatives of the profile.
Proposition~\ref{prop:wave-realization} collects these three conclusions
for one realization.

\subsection{Setting and the critical-point hypothesis}

Fix the coordinate cube $\Omega$ and the smooth coordinate extensions
used in Proposition~\ref{prop:spherical-profiles}.  Choose a sufficiently
small metric neighborhood $\mathcal G$ of the round metric inside
$\mathcal U$.  The coordinate and Riemannian norms on
$\overline\Omega$ are uniformly equivalent for $g\in\mathcal G$, and
the two-coordinate projections have the uniform Lipschitz bounds
for the ambient Riemannian distance established there.  We fix this neighborhood and these
coordinates for the rest of the argument.

In this section the incoming metric $g\in\mathcal G$ is smooth and equals
$g_*$ for $r\ge R_0$.  Fix radii
\[
 \frac15\le R_0<t<t_1<t_2<\frac3{10}.
\]
The fixed smooth profile $\phi$ equals
\[
 \phi_0=\log t+r^2-t^2
\]
off a compact subset of $\Omega$, has no critical point in
$\overline\Omega$, and satisfies
\begin{equation}
 \phi>\log r\quad(0<r\le R_0).
 \label{eq:wave-envelope-margin}
\end{equation}
Put $a=\grad_g\phi$ and $L=(1+|a|_g^2)^{1/2}$.  On
$K_1=\{r\le t_1\}$ we impose the following pointwise hypotheses:
\begin{enumerate}
\item At a noncritical point there is a vector $b$ such that
\begin{equation}
 b\perp a,\qquad |b|_g=L,\qquad
 \Hess_g\phi(a,a)+\Hess_g\phi(b,b)>0.
 \label{eq:wave-strict-hypothesis}
\end{equation}
\item At a critical point, $\Hess_g\phi$ is positive definite on some
two-dimensional subspace of the tangent space.
\end{enumerate}
The second hypothesis is stronger than the existential inequality at
$a=0$.  It is available for the profiles used here: all their critical
points in $K_1$ lie on $r=0$, where the perturbations vanish and
\[
 \Hess_g\phi=2(dx_1^2+dx_2^2).
\]
Here $d\phi=0$, so this Hessian identity is independent of the connection.
The form is positive semidefinite of rank two.  At every other point
outside the perturbations with $r<1$, $d(r^2)\ne0$.  Thus the
conclusions of Proposition~\ref{prop:spherical-profiles} supply exactly
the hypotheses above.

All profile data, the incoming metric, and the radii are fixed before
the integer $n$ is chosen.  Constants in phase estimates may depend on
these fixed data and on finitely many requested derivative orders.  The
constants in the final sign and area estimates will depend only on the
fixed coordinate and metric bounds.

\subsection{Complex Hessians below the profile}

Complex tangent vectors and symmetric forms are paired bilinearly in
the phase equations.  Transpose therefore means ordinary transpose,
without complex conjugation.

\begin{lemma}[Admissible phase Hessians]
\label{lem:phase-hessian}
Let $H$ be a real symmetric form on a three-dimensional Euclidean
space, and let $z=a+ib$, where $a,b$ are real,
$a\perp b$, and $|b|^2=1+|a|^2$.
If $a\ne0$ and $H(a,a)+H(b,b)>0$, there is a complex symmetric form
$Q$ such that
\begin{equation}
 Qz=0,\qquad \Re Q<H.
 \label{eq:phase-hessian-conditions}
\end{equation}
If $a=0$, the same conclusion follows from $H(b,b)>0$.
Moreover, a choice satisfying \eqref{eq:phase-hessian-conditions}
continues with locally bounded coefficients and a positive real-part
gap under small changes of $H,z$ with $z\cdot z=-1$.
\end{lemma}

\begin{proof}
Write $Q=U+iD$ with $U,D$ real symmetric.  The equation $Qz=0$ is
equivalent to
\[
 Db=Ua,\qquad Da=-Ub.
\]
For $a\ne0$, these prescribed columns of a symmetric $D$ are compatible
exactly when $U(a,a)+U(b,b)=0$.  To see sufficiency explicitly, write
$a=Ae_1$, $b=Be_2$ in an orthonormal frame.  Prescribe
\[
 De_1=-\frac BA Ue_2,\qquad De_2=\frac AB Ue_1.
\]
Their overlapping off-diagonal entries agree precisely when
$A^2U_{11}+B^2U_{22}=0$; symmetry prescribes the remaining entries in
the first two rows, and $D_{33}$ is free.  Taking
\[
 \alpha=\frac{H(a,a)+H(b,b)}{|a|^2+|b|^2}>0,
 \qquad U=H-\alpha\Id
\]
therefore proves the assertion.

If $a=0$, then $|b|=1$.  Set $U=0$ on $\R b$ and
$U=-M\Id$ on $b^\perp$, with no mixed entries, and take $D=0$.
In a frame starting with $b$, the first diagonal entry of $H-U$ is
$H(b,b)>0$.  Increasing $M$ makes its transverse Schur complement
positive definite.  Hence $H-U>0$ and $Qb=0$.

For the continuation assertion, start with $Qz=0$ and $z\cdot z=-1$.
For nearby $z'$ with $z'\cdot z'=-1$, define
\[
 P(z')=\Id+z'(z')^T,
 \qquad Q'=P(z')^TQP(z').
\]
Then $P(z')z'=0$, so $Q'$ is symmetric and annihilates $z'$.
At $z'=z$ one has $QP(z)=P(z)Q=Q$, and hence $Q'=Q$.
Although $P(z')$ is singular, these identities suffice: continuity
preserves the strict inequality $H'-\Re Q'>0$ for nearby $H',z'$.
No division by $|a|$ is involved in this continuation.
\end{proof}

For each $p\in K_1$, choose a unit vector $\nu$ with
$a(p)=|a(p)|\nu$; at a critical point any unit $\nu$ is allowed.
There are two linearly independent vectors $b_0,b_1\in\nu^\perp$,
both of length $L(p)$, for which the strict inequality in
\eqref{eq:wave-strict-hypothesis} holds.  At a noncritical point,
the good directions form a nonempty open subset of the transverse
circle.  At a critical point, every two-plane intersects the positive
two-plane of the Hessian in a nonzero positive direction.  Again the
good set in the circle of $\nu^\perp$ is nonempty and open.

These choices have uniform bounds for the fixed profile, including near
its critical points.  Indeed, consider the compact set
\[
 \mathcal K=\{(p,\nu):p\in K_1,\ |\nu|_g=1,
                         \ a(p)=|a(p)|_g\nu\}.
\]
Use finitely many smooth local orthonormal frames.  Near each datum in
$\mathcal K$, project and normalize the two chosen directions in the
nearby planes $\nu^\perp$.  Independence and strictness persist on a
sufficiently small neighborhood.  Lemma~\ref{lem:phase-hessian} continues
the two Hessians there.  It also handles a third vector with slightly
shifted real part.  A finite covering of $\mathcal K$ consequently gives
numbers $\delta_0,\kappa,\eta>0$ and bounded choices for the following
three data, for every $0\le\delta\le\delta_0$:
\begin{equation}
 \begin{split}
 z_0&=a+ib_0,\qquad z_1=a+ib_1,\\
 z_2&=a+\delta\nu+ib'_0,
 \qquad
 b'_0=\frac{(1+|a+\delta\nu|_g^2)^{1/2}}{L(p)}b_0.
 \end{split}
 \label{eq:three-phase-data}
\end{equation}
The independent imaginary directions $b_0,b_1$ will control derivatives
in $\nu^\perp$ in a real superposition.  The shift $\delta\nu$ will
separate its value from its derivative in the $\nu$ direction.
Lemma~\ref{lem:wave-noncancellation} proves the resulting real-jet estimate.
The associated Hessians satisfy
\begin{equation}
 Q_\ell z_\ell=0,\qquad
 \Re Q_\ell\le\Hess_g\phi(p)-4\kappa g(p),
 \qquad \ell=0,1,2,
 \label{eq:uniform-phase-gap}
\end{equation}
and the sine of the angle between $b_0,b_1$ is at least $\eta$.
All three vectors have bilinear square $-1$.  No globally continuous
selection of the choices on $\mathcal K$ is required.  The finite-cover
construction bounds each selected choice and its continuation uniformly;
these bounds, $\kappa$, and $\eta$ may depend on the fixed profile.

\subsection{A finite Taylor construction}

We use a finite Taylor form of the complex-phase WKB construction;
compare \cite[Section~3]{DenckerSjostrandZworski2004}.
The explicit truncation keeps track of the weighted residual and
derivative orders required later.

We first record the elementary algebra behind the formal equations.
Let $D_z=z\cdot\partial$ for a nonzero $z\in\C^3$, and let
$\mathcal P_d$ denote the complex homogeneous polynomials of degree $d$.
Then
\begin{equation}
 D_z:\mathcal P_{d+1}\longrightarrow\mathcal P_d
 \quad\hbox{is onto}.
 \label{eq:polynomial-surjectivity}
\end{equation}
For an explicit right inverse, put
\[
 \ell_z(x)=\frac{\overline z\cdot x}{\sum_i|z_i|^2},
 \qquad \pi_zx=x-\ell_z(x)z,
 \qquad
 (\mathcal I_z f)(x)=\ell_z(x)
      \int_0^1 f\bigl(\pi_zx+s\ell_z(x)z\bigr)\,ds.
\]
The integral is an integral of polynomial coefficients; its arguments
are evaluated in the polynomial's complex extension.
Since $D_z\ell_z=1$ and $D_z(\pi_zx)=0$, differentiating the polynomial
antiderivative gives $D_z\mathcal I_z f=f$.  The output is homogeneous
of degree $d+1$.  For each fixed degree, its coefficients are bounded
on every compact set of nonzero $z$.  This proves both surjectivity and
the local uniformity needed below.

\begin{lemma}[Localized finite-order waves]
\label{lem:finite-waves}
For every pair of nonnegative integers $m,D$, there are common normal
coordinate balls about the points of $K_1$ and three waves at every
center $p\in K_1$, for all sufficiently large integers $n$, with the
following properties.  Each wave is a smooth complex function on the
sphere supported in $r<t_2$ and has the form
\begin{equation}
 Z=\zeta e^{nS}V,
 \qquad V=\sum_{j=0}^J n^{-j}V_j,
 \label{eq:finite-wave-form}
\end{equation}
in its center's normal coordinates.  Here $S,V_j$ are complex
polynomials and $\zeta$ is a smooth cutoff equal to one on a common
smaller ball.  With $\delta=n^{-1/2}$, the three phases have initial
data
\[
 S(p)=\phi(p),\qquad \grad S(p)=z_\ell,
 \qquad \Hess S(p)=Q_\ell,
\]
as in \eqref{eq:three-phase-data}--\eqref{eq:uniform-phase-gap}.
Their amplitudes satisfy $V_0(p)=1$ and $V_j(p)=0$ for $j>0$.
On their supports,
\begin{equation}
 \Re S(x)-\phi(x)
 \le n^{-1/2}\dist_g(p,x)-c\dist_g(p,x)^2,
 \label{eq:wave-localization}
\end{equation}
with a common $c>0$.  Setting $\Lambda_n=n(n+2)$, one has
\begin{equation}
 \bigl|\nabla^j(\Delta_g+\Lambda_n)Z(x)\bigr|_g
       \le Cn^{-D}e^{n\phi(x)},\qquad 0\le j\le m.
 \label{eq:wave-individual-residual}
\end{equation}
For any further fixed finite collection of derivative orders, the
construction can also be required to satisfy
\begin{equation}
 |\nabla^jZ(x)|_g\le C_j n^j e^{n\phi(x)}
 \label{eq:wave-individual-derivatives}
\end{equation}
at those orders.  All the constants, polynomial degrees, and ball
radii are fixed before $n$; they may depend on $g,\phi,m,D$ and the
further derivative orders.
\end{lemma}

\begin{proof}
Use normal coordinates based on the local orthonormal frames above.
The fixed smooth metric has uniform coordinate derivative bounds of
every specified finite order on a common ball.  Put $z=z_\ell$ and
$Q=Q_\ell$.  The quadratic initial Taylor polynomial is
\[
 S_{\le2}(x)=\phi(p)+z\cdot x+\tfrac12x^TQx.
\]
The eikonal error $E=1+g^{ij}S_iS_j$ has zero constant term because
$z\cdot z=-1$, and zero linear term because the metric's first
derivatives vanish at the normal-coordinate origin and $Qz=0$.

Choose integers
\[
 J=D+m+1,\qquad K=2D+3m+6,
 \qquad N_j=K+J-j,\qquad L_S=K+J+1.
\]
Construct $S$ through degree $L_S$ by cancelling the eikonal
coefficients through degree $L_S-1$.  At degree $d\ge2$, the new
homogeneous term $S_{d+1}$ enters only as $2D_zS_{d+1}$; all other
terms are already known.  Equation~\eqref{eq:polynomial-surjectivity}
therefore solves each successive equation, preserving the given
quadratic data.  The resulting actual smooth eikonal error satisfies
\[
 E=O(|x|^{L_S}).
\]
Here and below these are Taylor-vanishing statements for the actual
smooth metric coefficients, not only identities between formal symbols.

Define
\[
 T=2\grad S\cdot\nabla+\Delta_gS+2,
\]
where the last two terms act by multiplication.  Construct $V_j$
through degree $N_j$, with the prescribed constant terms, successively
for $j=0,\ldots,J$.  Cancel, through degree $N_j-1$, the coefficients of
\[
 R_0=TV_0,\qquad
 R_j=TV_j+\Delta_gV_{j-1}\quad(j\ge1).
\]
At degree $d$ the new term of $V_j$ again enters as
$2D_z V_{j,d+1}$, so the same right inverse applies.  This finite
schedule supplies all forcing terms: the largest phase degree needed
is $N_0+1=L_S$, and for $j\ge1$ the coefficient of
$\Delta_gV_{j-1}$ of degree $d\le N_j-1$ needs $V_{j-1}$ only through
degree $d+2\le N_j+1=N_{j-1}$.  Consequently
\[
 R_j=O(|x|^{N_j})=O(|x|^K).
\]
Every operation uses finitely many metric coefficients and bounded
right inverses on the compact family of phase data.  Thus all the
constructed coefficients have uniform bounds for these fixed orders,
including as $\delta=n^{-1/2}$ tends to zero.

The Taylor expansion of $\Re S-\phi$ has linear part either zero or
$\delta\nu\cdot x$.  Equation~\eqref{eq:uniform-phase-gap} and a common
bound on third derivatives give, after decreasing the common ball,
\[
 \Re S(x)-\phi(x)\le\delta|x|-\kappa|x|^2.
\]
Normal-coordinate radius equals $\dist_g(p,x)$ on this ball, proving
\eqref{eq:wave-localization}.  Choose a radial cutoff supported
in this ball, equal to one on a common smaller ball, and make the
support small enough to lie in $r<t_2$ for every center.

Before applying the cutoff, direct differentiation gives
\begin{equation}
 (\Delta_g+\Lambda_n)(e^{nS}V)
 =e^{nS}\left(n^2EV+
       \sum_{j=0}^J n^{1-j}R_j+n^{-J}\Delta_gV_J\right).
 \label{eq:finite-wave-residual-expansion}
\end{equation}
For any fixed $q\ge0$,
\begin{equation}
 \sup_{r\ge0} r^q e^{\sqrt n r-cn r^2}
       \le C_{q,c}n^{-q/2};
 \label{eq:gaussian-envelope-moment}
\end{equation}
this follows by substituting $s=\sqrt n r$.  Taylor's Theorem bounds
the derivatives of a remainder vanishing to order $K$ by
$C|x|^{K-j}$ through order $j\le m$.  Differentiating the exponential
at most $m$ times costs at most $Cn^m$.  Thus the terms in
\eqref{eq:finite-wave-residual-expansion} involving Taylor
remainders, after at most $m$ derivatives and then division by
$e^{n\phi}$, are bounded by
\[
 C n^{2+m-(K-m)/2}=C n^{-D-1}.
\]
The last amplitude term is bounded by $Cn^{m-J}=Cn^{-D-1}$.
On the support of a derivative of the cutoff, the distance from the
center has a fixed positive lower bound.  For large $n$,
Equation~\eqref{eq:wave-localization} then gives a fixed negative gap
between $\Re S$ and $\phi$, so all cutoff errors are exponentially
small relative to $e^{n\phi}$.  This proves
\eqref{eq:wave-individual-residual}.  Direct differentiation
and \eqref{eq:gaussian-envelope-moment} with $q=0$ prove
\eqref{eq:wave-individual-derivatives}.  Additional fixed
derivative orders only require additional finite coefficient bounds.
\end{proof}

The local waves are now available with any prescribed finite accuracy.
We next combine them into one real function.  The combination must have
both a nonvanishing value--gradient pair and many strict sign changes.
We first define the sign functional that will measure the latter property;
its geometric and continuity properties do not depend on the wave
construction.

\subsection{A sign functional bounded by nodal area}

Subdivide $\Omega$, up to its faces, into a finite collection
$\mathscr P$ of open Euclidean cubes $P$ so fine that numbers $\ell_P$
can be chosen with
\begin{equation}
 \ell_P/2\le L(x)\le\ell_P\qquad(x\in P).
 \label{eq:sign-partition-scale}
\end{equation}
For example, use maxima of $L$ on the closed cubes of a sufficiently
fine regular subdivision; positivity and uniform continuity of $L$
ensure \eqref{eq:sign-partition-scale}.  Fix a sufficiently
small $\epsilon>0$ using only the common coordinate and metric
bounds, as specified below, and set $d_P=\epsilon/(n\ell_P)$.
For a continuous real function $f$ on the sphere define
\begin{equation}
 \mathcal F_n(f)=
 \sum_{P\in\mathscr P}\sum_{i=1}^3d_P^{-1}
 \int_{\{x\in P:x+d_Pe_i\in P\}}
       \mathbf1_{\{f(x)f(x+d_Pe_i)<0\}}\,dx.
 \label{eq:sign-functional}
\end{equation}
All shifts and integrals in this formula use the fixed coordinates
on $\Omega$.  The profile, subdivision, and $n$ are held fixed when
the functional is applied to other functions or metrics.

\begin{proposition}[Persistence and area comparison for sign tests]
\label{prop:sign-functional}
\label{lem:sign-area}
The functional $\mathcal F_n$ is lower semicontinuous for uniform
convergence of real functions.  It is unchanged by multiplication
by a nonzero real constant or by an everywhere positive continuous
function.  There is a constant $C_{\rm area}$, depending only on the
fixed coordinates and metric neighborhood, such that for every
continuous real $f$ and every $h\in\mathcal G$,
\begin{equation}
 \mathcal F_n(f)\le
       C_{\rm area}\HH_h^2(\{f=0\}).
 \label{eq:sign-area-bound}
\end{equation}
In particular $C_{\rm area}$ is independent of $\phi,\mathscr P$, and
$n$.  There is also a constant $C_{\rm sign}$ with the same
independence such that
\begin{equation}
 0\le\mathcal F_n(f)\le
       C_{\rm sign}n\int_\Omega L(x)\,dx.
 \label{eq:sign-deterministic-upper}
\end{equation}
\end{proposition}

\begin{proof}
If $f_j\to f$ uniformly, every strict opposite-sign test for $f$
eventually remains a strict opposite-sign test for $f_j$.  Fatou's
Lemma, applied to each integral in the finite sum, proves lower
semicontinuity.  The two invariances follow directly from the sign
of the product at its endpoints.

For the area estimate fix a cube $P$, a direction $i$, and write its
coordinate in that direction as $t$.  Slice at levels $t=s+kd_P$,
where $0\le s<d_P$ and $k\in\mathbb Z$.  For two consecutive levels
inside $P$, each transverse point with opposite endpoint signs has
an interior zero in the intervening open slab.  Hence the set of such
transverse points is contained in the coordinate projection of the
zeros in that slab.  By the uniform Lipschitz projection bound, its
two-dimensional Lebesgue measure is at most a common constant times
the $h$-Hausdorff measure of those zeros.  These sets are measurable:
the zeros in an open slab are a countable union of compact sets, and
so are their continuous projections.  For fixed $s$ the slabs are
disjoint.  Summing over $k$, integrating over $0\le s<d_P$, and
dividing by $d_P$ gives
\[
 d_P^{-1}\int_{\{x\in P:x+d_Pe_i\in P\}}
  \mathbf1_{\{f(x)f(x+d_Pe_i)<0\}}\,dx
       \le C\HH_h^2(\{f=0\}\cap P).
\]
This remains valid if the measure on the right is infinite.
The open cubes are disjoint, so summing over them and over three
directions proves \eqref{eq:sign-area-bound}.  No regularity
of the zero set has been used.

Finally, bounding each indicator by one gives
\[
 \mathcal F_n(f)\le\frac{3n}{\epsilon}
        \sum_P\ell_P|P|
       \le\frac{6n}{\epsilon}\int_\Omega L(x)\,dx,
\]
by \eqref{eq:sign-partition-scale}.
\end{proof}

\subsection{A random superposition with a nonvanishing real jet}

Choose a set of centers in $K_1$ of cardinality $O(n^3)$ so that every
point of $K_1$ is within distance $1/n$ of a center.  Such a set is
obtained by taking a maximal $1/n$-separated subset and using the
uniform lower volume bound for sufficiently small metric balls.  At
each center take the three waves of Lemma~\ref{lem:finite-waves}.
Assign independent complex coefficients $\gamma_Z$ whose real and
imaginary parts are independent standard real Gaussians, and put
\begin{equation}
 U_n=w_n+\sum_Z\Re(\gamma_Z Z),
 \qquad w_n=\Re(x_1+ix_2)^n,
 \qquad W_n=e^{n\phi}+r^n.
 \label{eq:wave-weight}
\end{equation}
The weight $W_n$ is positive and is $C^1$ for $n\ge2$; we do not
require it to be smooth at $r=0$ for odd $n$.

The function $w_n$ is the restriction of a homogeneous harmonic
polynomial of degree $n$ on $\R^4$.  The Euclidean Laplacian in polar
coordinates therefore gives $(\Delta_{g_*}+\Lambda_n)w_n=0$.
For each fixed $j$, its derivatives are bounded by
$C_jn^j r^{n-j}$ for large $n$, up to harmless lower-degree derivative
terms with the same bound on the sphere.  Away from $r=0$, the loss
of a fixed power of $r$ is bounded.  Near $r=0$, choose a fixed small
neighborhood on which $re^{-\phi}<1$ uniformly; then
\[
 r^{n-j}e^{-n\phi}
   =(re^{-\phi})^{n-j}e^{-j\phi}
\]
is bounded and is exponentially small there for fixed $j$.
These observations prove
\begin{equation}
 |\nabla^jw_n|_g\le C_jn^jW_n,
 \qquad |\nabla W_n|_g\le CnW_n.
 \label{eq:background-weight-bounds}
\end{equation}
Furthermore, $(\Delta_g+\Lambda_n)w_n$ is supported in $r\le R_0$.
Equation~\eqref{eq:wave-envelope-margin} gives
$\max_{r\le R_0}re^{-\phi}<1$, so the same argument bounds any fixed
number of its derivatives by any prescribed inverse power of $n$
times $e^{n\phi}$, for all sufficiently large $n$.

\begin{lemma}[Uniform exponent in real-jet noncancellation]
\label{lem:wave-noncancellation}
Fix any nonnegative integers $m_*,D_*$.  Construct the waves with
those residual requirements and with all derivative bounds needed
through order $\max(m_*,2)$.  There is a fixed exponent $B=110$,
independent of the profile and these orders, such that with probability
tending to one as $n\to\infty$,
\begin{equation}
 |U_n|+n^{-1}|\nabla U_n|_g\ge n^{-B}W_n
       \quad\hbox{everywhere on the sphere},
 \label{eq:wave-jet-lower}
\end{equation}
and simultaneously
\begin{equation}
 \begin{split}
 |\nabla^jU_n|_g&\le C_j n^{j+4}W_n,\\
 |\nabla^j(\Delta_g+\Lambda_n)U_n|_g
      &\le C n^{4-D_*}W_n,
                 \qquad 0\le j\le m_*.
 \end{split}
 \label{eq:wave-global-estimates}
\end{equation}
The constants and the rate of convergence of the probability may
depend on the fixed metric, profile, and orders, but not on $n$.
\end{lemma}

\begin{proof}
Let $G_n^{\rm coef}$ be the event that every coefficient has magnitude
at most $n$.  There are $O(n^3)$ coefficients, so Gaussian tails imply
$\Pr((G_n^{\rm coef})^c)\to0$.  On this event, summing the individual
wave estimates costs at most a factor $Cn^4$.  Together with
\eqref{eq:background-weight-bounds} and the background residual
estimate, this proves \eqref{eq:wave-global-estimates}.
In finitely many smooth local orthonormal frames, the normalized real
four-component jet
\[
 \mathcal J_n=\frac{(U_n,n^{-1}\nabla U_n)}{W_n}
\]
has Lipschitz constant at most $Cn^6$ on $G_n^{\rm coef}$.  Indeed,
its derivative involves only $U_n,\nabla U_n,\nabla^2U_n$,
$W_n^{-1}\nabla W_n$, and bounded derivatives of the frames;
the preceding estimates even give $Cn^5$.

First consider points where $r^n>n^6e^{n\phi}$.  The ambient rotation
field $R=x_1\partial_{x_2}-x_2\partial_{x_1}$ is tangent to the sphere
and has uniformly bounded $g$-length.  Writing $x_1+ix_2=re^{i\theta}$,
\[
 w_n=r^n\cos(n\theta),\qquad
 Rw_n=-nr^n\sin(n\theta).
\]
Consequently $|w_n|+n^{-1}|\nabla w_n|_g\ge c r^n$.  The wave jet
on $G_n^{\rm coef}$ is bounded by $Cn^4e^{n\phi}$, so domination
proves \eqref{eq:wave-jet-lower} in this region for large $n$.

The remaining closed set
\[
 E_n=\{r^n\le n^6e^{n\phi}\}
\]
lies in $K_1$ for large $n$.  Indeed $\phi=\phi_0$ when $r\ge t_1$,
and $\phi_0-\log r$ has a strictly negative maximum on that compact
region.  Fix $x\in E_n$ and take a center $p$ at distance at most
$1/n$.  The three associated cutoffs are one at $x$.  Their fixed
Taylor bounds and constant terms imply
\[
 V(x)=1+O(n^{-1}),\qquad
 \frac{\nabla Z(x)}{nZ(x)}=z_\ell+O(n^{-1}),\qquad
 \frac{|Z(x)|}{W_n(x)}\ge c n^{-6},
\]
in a common orthonormal frame, with the initial vectors at $p$
identified with that frame up to an $O(n^{-1})$ change.
For the last estimate, Taylor's Theorem gives
$n|\Re S(x)-\phi(x)|\le C$ at distance $1/n$, and
$W_n(x)\le(1+n^6)e^{n\phi(x)}$.

After factoring a complex phase and a positive modulus from each
column, the covariance supplied just by these three independent
coefficients is described, up to $O(n^{-1})$ column errors, by the
real and imaginary parts of
\[
 (1,a+ib_0),\qquad (1,a+ib_1),\qquad
 (1,a+\delta\nu+ib'_0),\qquad \delta=n^{-1/2}.
\]
The resulting real $4$ by $6$ matrix has least singular value at least
$c\delta$.  To check this, test its transpose on $(s,v)\in\R\times\R^3$.
The two transverse imaginary columns control the component of $v$
in $\nu^\perp$, using the positive angle bound and $L\ge1$.
Subtracting the first and third real-column evaluations gives
$\delta\langle v,\nu\rangle$.  The first real column then controls
$s$, since $a=|a|\nu$ and $|a|$ is bounded for the fixed profile.
Thus $|(s,v)|$ is bounded by $C\delta^{-1}$ times the transpose
evaluation norm.  The $O(n^{-1})$ errors preserve this estimate for
large $n$.

Multiplication of a complex coefficient by a unit complex number
preserves its Gaussian law.  Moreover, each factored modulus is at
least $cn^{-6}$.  Therefore every unit real linear functional of
$\mathcal J_n(x)$ has standard deviation at least $cn^{-13/2}$.
Adding the other independent waves cannot decrease covariance.
The deterministic jet of $w_n$ changes only the mean.  The Gaussian
density in four real dimensions is consequently at most $Cn^{26}$,
and hence at most $Cn^{28}$.  In particular,
\[
 \Pr\{ |\mathcal J_n(x)|\le2n^{-B}\}
       \le Cn^{28-4B},\qquad x\in E_n.
\]

In each of the finitely many frame charts take a net of $E_n$ with
points in $E_n$ and mesh at most $n^{-B-10}$.  Metric packing bounds
its total cardinality by $Cn^{3(B+10)}$.  A union bound makes the
probability of a small jet at any net point at most
\[
 Cn^{28-4B+3(B+10)}=Cn^{58-B}.
\]
For $B=110$ this tends to zero.  On $G_n^{\rm coef}$ the Lipschitz
error from a net point to a general point is $O(n^{-B-4})$, so the
jet norm is at least $n^{-B}$ throughout $E_n$ on the resulting
event.  Its Euclidean norm is bounded above by
$(|U_n|+n^{-1}|\nabla U_n|_g)/W_n$.  This proves the desired lower
bound everywhere.  The polynomial powers used in this argument do
not depend on the number of fixed Taylor or transport orders:
increasing those orders changes only constants, radii, and thresholds.
\end{proof}

\subsection{Expected size of the sign witness}

\begin{lemma}[Profile-independent sign gain]
\label{lem:wave-sign-expectation}
With a fixed sufficiently small $\epsilon>0$ as above, there is a
constant $c_{\rm sign}>0$, independent of the fixed profile and of
the wave accuracy orders, such that for all sufficiently large $n$,
\begin{equation}
 \mathbb E\mathcal F_n(U_n)
       \ge c_{\rm sign} n\int_\Omega L(x)\,dx.
 \label{eq:sign-expectation-lower}
\end{equation}
The threshold for $n$ may depend on the entire fixed construction.
\end{lemma}

\begin{proof}
Write the random part of $U_n$ as $X_n$ and let
\[
 \sigma_n(x)^2=\mathbb E X_n(x)^2=\sum_Z|Z(x)|^2.
\]
The wave at a center within distance $1/n$ shows that
\begin{equation}
 \sigma_n(x)\ge c e^{n\phi(x)},\qquad x\in\overline\Omega,
 \label{eq:sign-variance-lower}
\end{equation}
with a fixed-stage constant $c>0$.  Since
$\min_{\overline\Omega}(\phi-\log r)>0$, the standardized mean
$w_n(x)/\sigma_n(x)$ tends uniformly to zero on $\overline\Omega$.

Fix $P$ and $x\in P$ such that all three points
$y_i=x+d_Pe_i$ lie in $P$.  Use the same near-center set for every
axis: retain the waves whose centers have distance at most
$n^{-1/4}$ from $x$.  A center outside this set has distance at
least $\tfrac12n^{-1/4}$ from every $y_i$, because
$\dist_g(x,y_i)\le C d_P\le C\epsilon/n$.
Equation~\eqref{eq:wave-localization} therefore
bounds all these far waves, at $x$ and all $y_i$, by the respective
$e^{n\phi}$ factors times $C\exp(-c\sqrt n)$ for large $n$.
After summing $O(n^3)$ waves, their squared sums are still negligible
relative to the lower variances in
\eqref{eq:sign-variance-lower}.  Their covariance contributions
are negligible by Cauchy--Schwarz as well.

For every retained wave, the cutoff is one at both endpoints, and
uniform Taylor estimates give
\[
 V(x)=1+o(1),\qquad
 \Re dS(x)=d\phi(x)+o(1),\qquad
 |\beta_Z|_g=L(x)+o(1),\quad \beta_Z=\Im dS(x).
\]
The same estimates hold in the small segments between the endpoints.
Indeed their distances from the center are $O(n^{-1/4})$, the phase
and amplitude derivatives have fixed bounds, and the shifted initial
real gradient differs by only $n^{-1/2}\nu$.  Every error here is
uniform in centers, $P,x$, and $i$ for the fixed profile and orders.
Taylor expansion over a segment of length $d_P$ yields
\begin{equation}
 \frac{Z(y_i)}{Z(x)}
   =\alpha_i e^{i\theta_{Zi}}(1+o(1)),
 \qquad
 \alpha_i=e^{nd_P\partial_i\phi(x)},\quad
 \theta_{Zi}=nd_P\beta_Z(e_i).
 \label{eq:sign-wave-ratio}
\end{equation}
For example, the phase Taylor remainder after multiplication by $n$
is $O(nd_P^2)=O(n^{-1})$.  The real factor $\alpha_i$ is common to
all retained waves for the given axis.  It and its reciprocal are
bounded by a constant independent of the profile, since
$nd_P|\partial_i\phi(x)|\le C\epsilon$.

The covariance of $X_n(x)$ and $X_n(y_i)$ is
\[
 \sum_Z\Re\bigl(\overline{Z(x)}Z(y_i)\bigr).
\]
Discarding the negligible far terms and using
\eqref{eq:sign-wave-ratio} in both this numerator and the
variance at $y_i$ proves that their correlation satisfies
\begin{equation}
 \rho_i=\sum_{Z\text{ near}}\omega_Z\cos\theta_{Zi}+o(1),
 \qquad
 \omega_Z=\frac{|Z(x)|^2}
                  {\sum_{Z'\text{ near}}|Z'(x)|^2}.
 \label{eq:sign-common-weights}
\end{equation}
The same nonnegative weights, summing to one, apply to all three
axes: the common positive factor $\alpha_i$ cancels on standardizing.

Choose constants $c_0,C_0>0$ for the fixed equivalence of coordinate
and metric covector norms, so that
\[
 c_0|\beta|_g^2\le\sum_{i=1}^3\beta(e_i)^2
                         \le C_0|\beta|_g^2
\]
on $\overline\Omega$ for all the incoming metrics under consideration.
Choose $\epsilon$ so small that $2\epsilon\sqrt{C_0}\le1$.
By \eqref{eq:sign-partition-scale} and $L\ge1$, for large $n$
every retained wave satisfies
\begin{equation}
 |\theta_{Zi}|\le1,\qquad
 \sum_{i=1}^3\theta_{Zi}^2\ge\frac{c_0\epsilon^2}{16}.
 \label{eq:sign-angular-gap}
\end{equation}
For the lower estimate it suffices to use
$|\beta_Z|_g\ge L/2$ and $\ell_P\le2L$.  These deliberately loose
constants emphasize their independence from the profile.

Since $1-\cos s\ge s^2/4$ for $|s|\le1$, averaging
Equation~\eqref{eq:sign-angular-gap} with the common weights gives,
for all sufficiently large $n$,
\[
 \sum_{i=1}^3(1-\rho_i)\ge\frac{c_0\epsilon^2}{128}.
\]
Also every $\rho_i\ge\cos(1)+o(1)>0$.  Consequently at least one
axis has its correlation in a fixed compact subinterval of $(-1,1)$,
bounded away from $1$ by $c_0\epsilon^2/384$.  This reasoning still
applies if a single carrier direction accounts for nearly all the
variance.

For completeness, the required opposite-sign probability bound can
be read directly from bivariate Gaussian densities.  If the
correlation belongs to a fixed compact subinterval of $(-1,1)$ and
the two standardized means have absolute value at most $1/4$, that
density has a positive common minimum on the rectangle
$[1,2]\times[-2,-1]$.  Its integral gives a fixed positive probability
of opposite signs.  The standardized means here tend uniformly to
zero, so there is $p_0>0$, depending only on the fixed angular gap,
such that
\[
 \sum_{i=1}^3
   \Pr\{U_n(x)U_n(x+d_Pe_i)<0\}\ge p_0
\]
whenever all three shifts stay in $P$, for sufficiently large $n$.
No measurable choice of a preferred axis is needed in this inequality.

The subdivision is finite and fixed before $n$.  Thus, for large $n$,
the subset of each $P$ on which all three shifts remain in $P$ has
volume at least $|P|/2$.  Integrating the last probability estimate
and using Tonelli's Theorem yields
\[
 \mathbb E\mathcal F_n(U_n)
 \ge\frac{p_0 n}{2\epsilon}\sum_P\ell_P|P|
 \ge\frac{p_0 n}{2\epsilon}\int_\Omega L(x)\,dx.
\]
This proves Lemma~\ref{lem:wave-sign-expectation} with a profile-independent constant.  All
limits used to reach its threshold were taken with the profile and
finite accuracy orders fixed.
\end{proof}

\begin{proposition}[Simultaneous wave realization]
\label{prop:wave-superposition}
\label{prop:wave-realization}
Under the hypotheses of this section, fix arbitrary nonnegative
integers $m_*,D_*$.  For every sufficiently large integer $n$ there
is a smooth real $U_n$ satisfying
\eqref{eq:wave-jet-lower} and
\eqref{eq:wave-global-estimates}, with $B=110$, and
\begin{equation}
 \mathcal F_n(U_n)\ge
       c_{\rm wave} n\int_\Omega L(x)\,dx,
 \label{eq:wave-sign-lower}
\end{equation}
where $c_{\rm wave}>0$ depends only on the fixed coordinates and
metric neighborhood.  Moreover,
\begin{equation}
 U_n=w_n\quad\hbox{and}\quad g=g_*
       \qquad\hbox{on }\{r\ge t_2\}.
 \label{eq:wave-clean-exterior}
\end{equation}
The thresholds and constants in the differential estimates may
depend on the fixed profile and accuracy orders; $B$ and
$c_{\rm wave}$ do not.  The sign functional is finite, fixed, and
has all the persistence and area-comparison properties of
Proposition~\ref{prop:sign-functional}.
\end{proposition}

\begin{proof}
Let $G_n$ be the simultaneous event of
Lemma~\ref{lem:wave-noncancellation}, and put
$A_n=n\int_\Omega L(x)\,dx$.  By
Proposition~\ref{prop:sign-functional} and
Lemma~\ref{lem:wave-sign-expectation},
\[
 0\le\mathcal F_n(U_n)\le C_{\rm sign}A_n,
 \qquad \mathbb E\mathcal F_n(U_n)\ge c_{\rm sign}A_n,
 \qquad \Pr(G_n^c)\longrightarrow0.
\]
It follows without an independence assumption that
\[
 \mathbb E\bigl[\mathcal F_n(U_n)\mathbf1_{G_n}\bigr]
 \ge\bigl(c_{\rm sign}-C_{\rm sign}\Pr(G_n^c)\bigr)A_n
 \ge\tfrac12c_{\rm sign}A_n
\]
for all sufficiently large $n$.  Some realization in $G_n$ therefore
has $\mathcal F_n(U_n)\ge(c_{\rm sign}/4)A_n$.  Select such a
realization and set $c_{\rm wave}=c_{\rm sign}/4$.
All waves are supported in $r<t_2$, and $g$ is round for $r\ge R_0$.
Since $R_0<t_2$, Equation~\eqref{eq:wave-clean-exterior} follows.
\end{proof}

\begin{remark}[Order of choices]
\label{rem:wave-parameter-order}
The profile is constructed first and is used here only after it has
been fixed.  Next one specifies the finite residual and derivative
orders required downstream.  The phase and amplitude degrees and
their common support radii are then fixed on the compact family with
$0\le\delta\le\delta_0$.  Finally $n$ is chosen so that
$\delta=n^{-1/2}$ and all preceding thresholds hold.  Neither a small
Hessian gap nor a poorly conditioned pair of transverse directions
changes the exponent $B$; it only enlarges these thresholds.  The
sign constant comes from \eqref{eq:sign-angular-gap}, which
uses $L/\ell_P$ and the fixed norm equivalence, not those
profile-dependent phase constants.  This distinction is what permits
the gradient integral to be made arbitrarily large before choosing
the frequency.
\end{remark}
